\par\noindent\begin{minipage}{\linewidth}
% !TeX root = ../教材题库_全部条目_转写工作稿.tex
% Source: XJ-B1-C01-PR-P021-Y0793
% 第1章 / 1.2 常用逻辑用语 / 教材 p.21 / PDF 第28页
% 已对照原页转写题干及公式；待教师复核；未制作答案
% 原题及相邻原解截图（核对用）：../题目原图/XJ-B1-C01-PR-P021-Y0793-01.webp
% 本文件只在转写完成后作为题干使用；OCR 条目尚不是可靠题干。
\begin{bankitem}{p.21 练习 第1题}{第1章 / 1.2 常用逻辑用语 / 教材 p.21 / PDF 第28页}{XJ-B1-C01-PR-P021-Y0793}{已对照原页转写题干及公式；待教师复核；未制作答案}
指出下列命题中使用了什么量词以及量词的作用范围，并把量词用相应的数学符号取代：
\subq{1}对任意整数 $n$，$2n+1$ 是奇数；
\subq{2}对某个有理数 $x$，有 $4x=\dfrac12$；
\subq{3}线段 $AB$ 上有一点 $M$ 满足比例式 $\dfrac{AM}{MB}=\dfrac{MB}{AB}$。
\end{bankitem}
\end{minipage}\par\medskip
